\subsection{PRINCIPAL ELEMENTS}

DEFINITION 3. Assume that g is semi-simple.

\begin{enumerate}
\def\labelenumi{(\roman{enumi})}
\item
  A nilpotent element \(x\) of \(\mathfrak{g}\) is said to be principal if the dimension of \(\operatorname{Ker} \operatorname{ad} x\) is the rank of \(\mathfrak{g}\) .
\item
  A simple element \(h\) of \(\mathfrak{g}\) is said to be principal if \(h\) is regular and the eigenvalues of \(\operatorname{ad} h\) in an algebraic closure of \(k\) belong to 2Z.
\item
  An \(\mathfrak{sl}_2\) -triplet \((x,h,y)\) of \(\mathfrak{g}\) is said to be principal if the length of \(\mathfrak{g}\) , considered as a module over \(kx + kh + ky\) , is equal to the rank of \(\mathfrak{g}\) .
\end{enumerate}

PROPOSITION 7. Assume that g is semi-simple. Let \((x,h,y)\) be an \(sl_{2}\) -triplet in g. The following conditions are equivalent:

\begin{enumerate}
\def\labelenumi{(\roman{enumi})}
\item
  \(x\) is principal;
\item
  h is principal;
\item
  \((x,h,y)\) is principal.
\end{enumerate}

For \(p \in Z\) , let \(\mathfrak{g}^{p} = \operatorname{Ker}(\operatorname{ad} h - p)\) . Let \(g' = \sum_{p \in Z} g^{2p}\) . If g is considered as a module over \(a = kx + kh + ky\) , \(g'\) is the sum of the simple submodules of odd dimension (§1, no. 2, Cor. of Prop. 2). Let l (resp. \(l'\) ) be the length of g (resp. \(g'\) ) considered as an a-module. By §1, no. 2,

\[
\dim (\operatorname{Ker} \operatorname{ad} x) = l \geq l ^ {\prime} = \dim (\operatorname{Ker} \operatorname{ad} h) \geq \operatorname{rk} (\mathfrak {g}).
\]

The equivalence of (i) and (iii) follows immediately. On the other hand, condition (ii) means that \(\dim(\operatorname{Ker} \operatorname{ad} h) = \operatorname{rk}(\mathfrak{g})\) and \(g' = g\) , in other words that

\[
\dim (\operatorname{Ker} \operatorname{ad} h) = \operatorname{rk} (\mathfrak {g})
\]

and \(l = l'\) . The equivalence of (ii) and the other conditions follows.

PROPOSITION 8. Assume that g is semi-simple \(\neq 0\) . Let h be a splitting Cartan subalgebra of g, R the root system of \((\mathfrak{g},\mathfrak{h})\) , B a basis of R, \(h^{0}\) the element of h such that \(\alpha(h^{0}) = 2\) for all \(\alpha \in B\) .

\begin{enumerate}
\def\labelenumi{(\roman{enumi})}
\item
  The element \(h^0\) is simple and principal.
\item
  The elements \(x\) of \(\mathfrak{g}\) such that there exists an \(\mathfrak{sl}_2\) -triplet of the form \((x, h^0, y)\) are the elements of \(\sum_{\alpha \in B} \mathfrak{g}^\alpha\) that have a non-zero component in each \(\mathfrak{g}^\alpha\) .
\end{enumerate}

The element \(h^0\) is that considered in §7, no. 5, Lemma 2 (cf.~loc. cit., formula (1)). It follows from this lemma that \(h^0\) is simple principal and that, if \(x \in \sum_{\alpha \in \mathrm{B}} \mathfrak{g}^\alpha\) has a non-zero component in each \(\mathfrak{g}^\alpha\) , there exists an \(\mathfrak{sl}_2\) -triplet of the form \((x, h^0, y)\) . Conversely, let \((x, h^0, y)\) be a \(\mathfrak{sl}_2\) -triplet. We have \([h^0, x] = 2x\) , so \(x \in \sum_{\gamma \in \mathrm{R}, \gamma(h^0) = 2} \mathfrak{g}^\gamma = \sum_{\alpha \in \mathrm{B}} \mathfrak{g}^\alpha\) . Similarly, \(y \in \sum_{\alpha \in \mathrm{B}} \mathfrak{g}^{-\alpha}\) . Write

\[
h ^ {0} = \sum_ {\alpha \in \mathrm{B}} a _ {\alpha} H _ {\alpha} \quad \text { where } a _ {\alpha} > 0 \text { for all } \alpha \in \mathrm{B},
\]

\[
x = \sum_ {\alpha \in \mathrm{B}} X _ {\alpha} \qquad \text { where } X _ {\alpha} \in \mathfrak {g} ^ {\alpha} \text { for   all } \alpha \in \mathrm{B},
\]

\[
y = \sum_ {\alpha \in \mathrm{B}} X _ {- \alpha} \quad \text { where } X _ {- \alpha} \in \mathfrak {g} ^ {- \alpha} \text { for   all } \alpha \in \mathrm{B}.
\]

Then

\[
\sum_ {\alpha \in \mathrm{B}} a _ {\alpha} H _ {\alpha} = h ^ {0} = [ y, x ] = \sum_ {\alpha , \beta \in \mathrm{B}} [ X _ {- \beta}, X _ {\alpha} ] = \sum_ {\alpha \in \mathrm{B}} [ X _ {- \alpha}, X _ {\alpha} ]
\]

so \([X_{-\alpha}, X_{\alpha}] \neq 0\) for all \(\alpha \in \mathrm{B}\) .

COROLLARY. In a splittable semi-simple Lie algebra, there exist principal nilpotent elements.

In a non-splittable semi-simple Lie algebra, 0 may be the only nilpotent element.

PROPOSITION 9. Assume that \(k\) is algebraically closed and that \(\mathfrak{g}\) is semi-simple. All the principal simple (resp. nilpotent) elements of \(\mathfrak{g}\) are conjugate under \(\operatorname{Aut}_e(\mathfrak{g})\) .

We retain the notations of Prop. 8. Let h be a principal simple element. It is conjugate under \(\operatorname{Aut}_{e}(\mathfrak{g})\) to an \(h' \in h\) such that \(\alpha(h') \in \{0,1,2\}\) for all \(\alpha \in B\) (no. 3, Prop. 5). Since \(h'\) is principal simple, \(\alpha(h') \neq 0\) and \(\alpha(h') \in 2Z\) for all \(\alpha \in B\) , so \(\alpha(h') = 2\) for all \(\alpha \in B\) , and hence \(h' = h^{0}\) . This proves the assertion for principal simple elements.

Let \(x, x'\) be principal nilpotent elements. There exist \(\mathfrak{sl}_2\) -triplets \((x, h, y)\) , \((x', h', y')\) . By Prop. 7, \(h\) and \(h'\) are principal simple, hence conjugate under \(\mathrm{Aut}_e(\mathfrak{g})\) by the preceding. So \(x\) and \(x'\) are conjugate under \(\mathrm{Aut}_e(\mathfrak{g})\) (Prop. 6).

Lemma 9. With the notations of Prop. 8, put \(\mathfrak{g}^p = \mathrm{Ker}(\mathrm{ad}h^0 - p)\) for \(p \in \mathbf{Z}\) . Let \(\mathfrak{g}_*^2\) be the set of elements of \(\mathfrak{g}^2 = \sum_{\alpha \in \mathrm{B}} \mathfrak{g}^\alpha\) that have a non-zero component in each \(\mathfrak{g}^\alpha\) . Let \(\mathrm{R}_+\) be the set of positive roots relative to \(\mathrm{B}\) , \(\mathfrak{n}_+ = \sum_{\alpha \in \mathrm{R}_+} \mathfrak{g}^\alpha\) , and \(x \in \mathfrak{g}_*^2\) . Then \(e^{\mathrm{ad}\mathfrak{n}_+}. x = x + [\mathfrak{n}_+, \mathfrak{n}_+]\) .

It is clear that \(e^{\mathrm{ad} \mathfrak{n}_+}.x \subset x + [\mathfrak{n}_+, \mathfrak{n}_+]\) . We prove that, if \(v \in [\mathfrak{n}_+, \mathfrak{n}_+]\) , then \(x + v \in e^{\mathrm{ad} \mathfrak{n}_+}.x\) . Put \(\mathfrak{n}^{(p)} = \sum_{r \geq p} \mathfrak{g}^{2r}\) ; it suffices to prove that

\[
x + v \in e ^ {\mathrm{ad} \mathfrak {n} _ {+}}. x + \mathfrak {n} ^ {(p)}
\]

for all \(p \geq 2\) . This is clear for \(p = 2\) since \(\mathfrak{n}^{(2)} = [\mathfrak{n}_+, \mathfrak{n}_+]\) (§3, no. 3, Prop. 9 (iii)). Assume that we have found \(z \in \mathfrak{n}_+\) such that \(v + x - e^{\mathrm{ad}z}.x \in \mathfrak{n}^{(p)}\) . Since there exists an \(\mathfrak{sl}_2\) -triplet of the form \((x, h^0, y)\) (Prop. 8), §1, no. 2, Cor. of Prop. 2 proves that \([x, \mathfrak{g}^{2p-2}] = \mathfrak{g}^{2p}\) ; hence, there exists \(z' \in \mathfrak{g}^{2p-2} \subset \mathfrak{n}_+\) such that

\[
v + x - e ^ {\mathrm{ad} z}. x \in [ z ^ {\prime}, x ] + \mathfrak {n} ^ {(p + 1)}.
\]

So \(v + x \in e^{\mathrm{ad}(z + z')}.x + \mathfrak{n}^{(p + 1)}\) , and our assertion is established by induction.

PROPOSITION 10. Assume that g is semi-simple. Let h be a splitting Cartan subalgebra of g, R the root system of \((\mathfrak{g},\mathfrak{h})\) , B a basis of R, \(R_{+}\) the set of positive roots relative to B, and \(n_{+}=\sum_{\alpha\in R_{+}}g^{\alpha}\) . The principal nilpotent elements belonging to \(n_{+}\) are the elements of \(n_{+}\) having a non-zero component in \(g^{\alpha}\) for all \(\alpha\in B\) .

Prop. 8 and Lemma 9 prove that such elements are principal nilpotent. We prove the converse. Evidently we can assume that \(\mathfrak{g}\) is simple. Let \(h^0\) and \(\mathfrak{g}^p\) be as in Prop. 8 and Lemma 9. Let \(\omega\) be the highest root, and put \(\omega(h^0) = 2q\) ; we have \(q = h - 1\) , where \(h\) is the Coxeter number of R, cf.~Chap. VI, §1, no. 11, Prop. 31. Then \(\mathfrak{g}^{2q} = \mathfrak{g}^\omega\) , \(\mathfrak{g}^{-2q} = \mathfrak{g}^{-\omega}\) , and \(\mathfrak{g}^{2k} = 0\) for \(|k| > q\) . There exists a principal \(\mathfrak{sl}_2\) -triplet \((x^0, h^0, y^0)\) . By §1, no. 2, Cor. of Prop. 2, \((\mathrm{ad} x^0)^{2q}(\mathfrak{g}^{-\omega}) = \mathfrak{g}^\omega\) , so \((\mathrm{ad} x^0)^{2q} \neq 0\) . Let \(x\) be a principal nilpotent element of \(\mathfrak{g}\) belonging to \(\mathfrak{n}_+\) . If \(\bar{k}\) is an algebraic closure of \(k\) , \(x \otimes 1\) and \(x^0 \otimes 1\) are conjugate under an automorphism of \(\mathfrak{g} \otimes_k \bar{k}\) (Prop. 9), so \((\mathrm{ad} x)^{2q} \neq 0\) . There exists \(\lambda \in \mathbb{R}\) such that \((\mathrm{ad} x)^{2q} \mathfrak{g}^\lambda \neq 0\) . Put \(x = \sum_{n \geq 1} x_n\) , where \(x_n \in \mathfrak{g}^{2n}\) . Then

\[
(\operatorname{ad} x) ^ {2 q} \mathfrak {g} ^ {\lambda} \subset (\operatorname{ad} x _ {1}) ^ {2 q} \mathfrak {g} ^ {\lambda} + \sum_ {k > 4 q + \lambda (h ^ {0})} \mathfrak {g} ^ {k} = (\operatorname{ad} x _ {1}) ^ {2 q} \mathfrak {g} ^ {\lambda},
\]

since \(4q + \lambda(h^0) \geq 4q - 2q = 2q\) . Now \((\operatorname{ad} x_1)^{2q} \mathfrak{g}^\lambda \subset \mathfrak{g}^{4q + \lambda(h^0)}\) , where \(\lambda = -\omega\) . Thus, \((\operatorname{ad} x_1)^{2q} \mathfrak{g}^{-\omega} = \mathfrak{g}^\omega\) . We have \(\omega = \sum_{\alpha \in \mathrm{B}} n_\alpha \alpha\) with \(n_\alpha > 0\) for all \(\alpha \in \mathrm{B}\) (Chap. VI, §1, no. 8, Remark). If there exists \(\alpha_0 \in \mathrm{B}\) such that \(x_1 \in \sum_{\alpha \in \mathrm{B}, \alpha \neq \alpha_0} \mathfrak{g}^\alpha\) , the relation

\[
\omega \notin - \omega + \sum_ {\alpha \in \mathrm{B}, \alpha \neq \alpha_ {0}} k \alpha
\]

implies that \(\mathfrak{g}^{\omega} \not\subset (\operatorname{ad} x_1)^p \mathfrak{g}^{-\omega}\) for all \(p\) ; this is absurd, so the component of \(x_1\) in \(\mathfrak{g}^{\alpha}\) is non-zero for all \(\alpha \in B\) .

